\documentclass[11pt,a4paper]{ctexart}

\usepackage[margin=2.2cm]{geometry}
\usepackage{amsmath,amssymb}
\usepackage{booktabs}
\usepackage{graphicx}
\usepackage{hyperref}
\usepackage{xcolor}
\usepackage{enumitem}

\hypersetup{colorlinks=true,linkcolor=blue!50!black,urlcolor=blue!50!black,citecolor=blue!50!black}
\newcommand{\safeincludegraphics}[2][]{%
  \IfFileExists{#2}{\includegraphics[#1]{#2}}{%
    \fbox{\begin{minipage}{0.42\linewidth}\centering
    缺少生成图\\\texttt{#2}\\请先运行\texttt{make data}
    \end{minipage}}%
  }%
}

\title{\heiti 钱学森问题的扩展求解\\
\large 含月球借力与发射窗口优化的绕日返回轨道设计}
\author{人工智能与数学软件期末项目}
\date{2026年}

\begin{document}
\maketitle

\begin{abstract}
本文在钱学森《星际航行概论》第5--6章拼接圆锥曲线理论基础上，完成了一个二维黄道面内的月球借力绕日返回轨道设计。程序实现了原始解析公式代码化、Sun--Earth--Moon--Rocket数值积分、JPL Horizons离线历表对照、月球双曲线借力解析与数值检验、固定发射日完整轨道求解、2026全年发射窗口扫描以及灵敏度分析。针对修正规则，主任务使用真实历表闭合：月球交会取交会时刻的真实月球位置，返回取返回时刻的真实地球位置，月球被动偏折不能闭合的速度矢量差全部计入$\Delta v_{\rm Moon,res}$。最优真实闭合结果给出发射日期为2026-08-12，近月距为1838.0 km，日心近日距为0.248 AU，总速度增量为20.099 km/s。
\end{abstract}

\section{问题与模型}

钱学森在《星际航行概论》中用拼接圆锥曲线法分析了从地球出发、绕过太阳后返回地球的轨道问题\cite{qian2008}。原算例只考虑日心二体椭圆和地球逃逸速度，未利用月球引力，也未对发射日期做全年扫描。本项目将问题写成
\begin{equation}
 \Delta v_{\rm total}=|\Delta v_{\rm launch}|+|\Delta v_{\rm Moon,res}|+|\Delta v_{\rm reentry}|.
\end{equation}
参考系取J2000黄道面二维投影，单位为km和s。太阳、地球、月球使用课程提供的2026年JPL Horizons离线缓存；为复核2026年发射后两年内返回，另用同一查询口径生成了2026--2029年Sun--Earth--Moon离线缓存，查询中心均为\texttt{@10}。火箭作为试探粒子。

为便于机器复核，程序额外输出任务摘要JSON、事件CSV和轨道CSV。这些文件包含发射时刻、月球交会时刻、返回时刻、近月距、近日距、三段速度增量、节能比例和带时间戳的轨道状态序列。JPL Horizons在线接口位于\texttt{src/jpl\_access.py}，其中\texttt{CENTER}参数固定为\texttt{@10}。

\section{M1：拼接圆锥曲线代码化}

设地球日心距离为$r_1$，日心转移椭圆近日距为$r_p$。当出发点为远日点时，
\begin{align}
 a&=\frac{r_1+r_p}{2},&
 e&=\frac{r_1-r_p}{r_1+r_p},&
 p&=\frac{2r_1r_p}{r_1+r_p}.
\end{align}
由活力公式
\begin{equation}
 v^2=\mu_\odot\left(\frac{2}{r}-\frac{1}{a}\right)
\end{equation}
得到远日点速度
\begin{equation}
 v_a=\sqrt{\frac{\mu_\odot}{r_1}\frac{2r_p}{r_1+r_p}},
 \qquad
 \Delta v_{\rm helio}=v_a-\sqrt{\mu_\odot/r_1}.
\end{equation}
地面发射速度按钱学森第一方案写为
\begin{equation}
 v_{\rm launch}=\sqrt{v_2^2+|\Delta v_{\rm helio}|^2}.
\end{equation}

代码位于\texttt{src/conics.py}。对$r_p=0.2$ AU的复核结果见表\ref{tab:m1}，最大相对误差为$8.02\times10^{-4}$，满足0.1\%要求。

\begin{table}[htbp]
\centering
\caption{M1：$r_p=0.2$ AU基准算例复核}
\label{tab:m1}
\begin{tabular}{lcc}
\toprule
量 & 程序结果 & 钱学森算例/原报告值\\
\midrule
$a$ / AU & 0.6000 & 0.6000\\
$e$ & 0.6667 & 0.6667\\
$p$ / AU & 0.3333 & 0.3333\\
$v_a$ / km s$^{-1}$ & 17.196 & 17.21\\
$\Delta v_{\rm helio}$ / km s$^{-1}$ & -12.588 & -12.59\\
$v_{\rm launch}$ / km s$^{-1}$ & 16.836 & 16.84\\
周期 / d & 169.756 & 169.8\\
\bottomrule
\end{tabular}
\end{table}

\section{M2--M3：N体积分器与历表对照}

\texttt{nbody.py}实现了加速度计算、单步推进和整段传播三类接口。积分器为二阶Velocity--Verlet：
\begin{align}
 r_{n+1}&=r_n+v_nh+\frac12a_nh^2,\\
 v_{n+1}&=v_n+\frac12(a_n+a_{n+1})h.
\end{align}

在无量纲二体圆轨道基准$\mu=4\pi^2$中，步长$1/4000$年积分1年后位置相对误差为$5.17\times10^{-6}$，优于$10^{-4}$。Sun--Earth--Moon三体从2026-01-01 Horizons状态出发，用$h=900$ s积分一年，并转回模拟太阳质心后与每日历表比较。最大位置残差为：Sun 0 km，Earth 5894 km，Moon 5918 km；三组相对位置残差为Moon--Earth 79.84 km、Moon--Sun 5917.60 km、Earth--Sun 5894.49 km，均满足6000 km要求。残差曲线见图\ref{fig:residual}。

\begin{figure}[htbp]
\centering
\safeincludegraphics[width=0.82\linewidth]{data/generated/horizons_residuals.png}
\caption{M3：Sun--Earth--Moon数值传播与JPL Horizons每日位置残差。}
\label{fig:residual}
\end{figure}

\section{M4：月球借力解析与数值实现}

月球作用范围内采用月心二体双曲线。若入射月心无穷远速度为$v_\infty$，近月距为$r_m$，双曲线偏折角为\cite{vallado2013,curtis2014}
\begin{equation}
 \delta=2\arcsin\frac{1}{1+r_m v_\infty^2/\mu_M}.
\end{equation}
程序在月心系将$v_{\infty,\rm in}$旋转$\pm\delta$得到$v_{\infty,\rm out}$，再加回月球日心速度。符号对应leading/trailing两侧绕飞。数值检验从近月点正反向积分至月球SOI，$v_\infty=1.2$ km/s、$r_m=5000$ km时，解析偏折角为47.794度，有限SOI数值偏折角为47.482度，差0.313度。

\section{M5--M6：单点轨道与全年扫描}

主任务代码位于\texttt{src/real\_closure.py}。对每个发射日期$t_0$，程序扫描地月转移时间、返回时间、$r_m\in[R_M+100,50000]$ km和两种绕飞侧。地月段用Lambert问题从真实地球位置连接到真实月球近月点；月球后日心段再用Lambert问题连接到返回时刻的真实地球位置。若月心双曲线最大偏折不能把$v_{\infty,\rm in}$转到所需$v_{\infty,\rm out}$，剩余速度矢量范数记入$\Delta v_{\rm Moon,res}$。旧的理想解析筛选模型仍保存在\texttt{legacy\_analytic\_scan.json}，只作对照，不作为主任务真实性复核依据。图\ref{fig:orbit}给出轨道示意，图\ref{fig:scan}给出全年扫描曲线和按实际近日距分箱的等值图。

\begin{figure}[htbp]
\centering
\safeincludegraphics[width=0.68\linewidth]{data/generated/orbit.png}
\caption{最优解对应的日心转移椭圆。}
\label{fig:orbit}
\end{figure}

\begin{figure}[htbp]
\centering
\safeincludegraphics[width=0.48\linewidth]{data/generated/scan_curve.png}
\safeincludegraphics[width=0.48\linewidth]{data/generated/scan_contour.png}
\caption{M6：2026全年发射窗口扫描曲线与$(t_0,r_p)$等值图。}
\label{fig:scan}
\end{figure}

最优结果见表\ref{tab:best}。固定最优发射日的无月球借力Lambert闭合基准为25.046 km/s，月球借力真实闭合方案为20.099 km/s，节省比例为19.8\%。该主解在2026-08-15到达真实月球附近，最近月距1838.0 km；随后进入日心转移，于2026-11-04经过近日点，并在2028-08-02 22:12 UTC与真实地球位置闭合，返回脱靶为0 km。

以2026-06-19作为输入发射日时，同一程序生成今日解摘要和轨道CSV。该日求得近月距10000.0 km、近日距0.245 AU，总速度增量20.992 km/s；无月借对照为25.262 km/s，节省比例16.9\%。该结果与全年扫描使用同一速度增量口径。

\begin{table}[htbp]
\centering
\caption{M5--M6：最优发射窗口与速度增量}
\label{tab:best}
\begin{tabular}{lc}
\toprule
项目 & 数值\\
\midrule
发射日期$t_0^\ast$ & 2026-08-12\\
月球交会时刻 & 2026-08-15\\
返回地球时刻 & 2028-08-02 22:12 UTC\\
近月距$r_m^\ast$ & 1838.0 km\\
绕飞侧 & leading\\
近日距$r_p^\ast$ & 0.248 AU\\
$\Delta v_{\rm launch}$ & 11.275 km/s\\
$\Delta v_{\rm Moon,res}$ & 8.824 km/s\\
$\Delta v_{\rm reentry}$ & 0.000 km/s\\
$\Delta v_{\rm total}$ & 20.099 km/s\\
返回$v_\infty$ & 11.124 km/s\\
总飞行时间 & 721.93 d\\
\bottomrule
\end{tabular}
\end{table}

\section{M7：灵敏度与守恒性}

图\ref{fig:energy}给出$h=3600$ s时一年传播的能量相对漂移，最大值为$8.64\times10^{-9}$，明显低于$10^{-6}$约束。图\ref{fig:sens}展示最优解附近的三类灵敏度。发射日期偏移$\pm5$天时总速度增量在20.10--21.31 km/s量级变化；近月距从1838 km增大到2297.5 km时总量从20.099 km/s升至20.208 km/s，反映偏折能力下降；步长从3600 s减小至900 s时一年三体能量相对漂移从$8.64\times10^{-9}$降至$5.40\times10^{-10}$。

\begin{figure}[htbp]
\centering
\safeincludegraphics[width=0.82\linewidth]{data/generated/energy_conservation.png}
\caption{M8：Sun--Earth--Moon一年传播的能量守恒监测。}
\label{fig:energy}
\end{figure}

\begin{figure}[htbp]
\centering
\safeincludegraphics[width=0.95\linewidth]{data/generated/sensitivity.png}
\caption{M7：发射日期、近月距和积分步长灵敏度。}
\label{fig:sens}
\end{figure}

\section{M8：可视化与可复现性}

所有图表由\texttt{run\_all.py}生成，包括轨道图、能量守恒图、扫描曲线、残差图和灵敏度图；动画输出为MP4文件。仓库中的\texttt{Makefile}提供一键生成、单独生成数据、单独编译报告和清理目标。关键数值保存在\texttt{results.json}，逐日最优值保存在CSV表中。

此外，\texttt{run\_all.py}会写出\texttt{grading\_check.json}，把M1基准误差、二体积分误差、三组Horizons相对残差、今日输入解、机器可读字段、真实月球交会、真实地球返回、速度预算闭合和365日扫描行数汇总到同一文件；也可单独运行\texttt{python src/self\_check.py}复核已生成文件。

\section{选做加分项 O5：轨道动画}

本文完成选做项O5“轨道动画”：\texttt{run\_all.py}会生成MP4动画文件，展示火箭沿最优日心转移椭圆绕日返回的过程。动画与静态轨道图使用同一组最优扫描参数，作为展示材料而不参与数值评分指标。

\section{选做项：3D、Lambert、多次借力、交互与创新扩展}

\subsection*{O1：3D月球倾角}

\texttt{src/extensions.py}读取Horizons缓存中的三维地月状态，保留月球相对地球的$z$分量，并用月球轨道标称倾角$5.145^\circ$校核二维黄道面模型。主最优日的地月面外距离为11731 km，瞬时倾角为1.84度。若把该面外量按一次平面修正估计，月心入射速度对应的速度惩罚约$9.0\times10^{-4}$ km/s。进一步对365个发射日逐日加上同一口径的3D惩罚后，二维真实闭合粗扫最优日2026-08-11的估计总量由20.268 km/s变为20.269 km/s，最优日期不移动。全年平均惩罚0.0029 km/s，最大惩罚0.0097 km/s，因此月球倾角对本模型的最优窗口影响小于主速度预算。

\subsection*{O2：广义相对论修正}

同一模块实现一阶Schwarzschild修正加速度
\[
\mathbf a_{\rm GR}=\frac{\mu_\odot}{c^2r^3}\left[\left(\frac{4\mu_\odot}{r}-v^2\right)\mathbf r
        +4(\mathbf r\cdot\mathbf v)\mathbf v\right].
\]
对主最优日心段，近日点速度为75.793 km/s，每轨道近日点进动为0.097角秒；近日点处牛顿加速度为$9.676\times10^{-5}$ km/s$^2$，GR修正加速度为$9.247\times10^{-12}$ km/s$^2$，相对量级$9.56\times10^{-8}$。这说明相对论效应可定量写入模型，但对本课程尺度的速度预算影响很小。

\subsection*{O3：自实现Lambert求解器}

\texttt{src/lambert.py}实现通用变量Lambert求解器，包含Stumpff函数、飞行时间方程和二分迭代，不调用\texttt{scipy.optimize}。单位圆四分之一周期验证中，求得首末速度范数均为1，速度误差为$3.7\times10^{-10}$。该求解器也被O4多次借力探索用于Moon--Venus和Venus--Earth两段拼接。

\subsection*{O4：Earth--Moon--Venus--Sun--Earth多次借力探索}

\texttt{src/multiflyby.py}探索了Earth $\to$ Moon $\to$ Venus $\to$ Sun $\to$ Earth序列。地球、月球使用真实闭合主解中的Horizons状态；金星使用本文新增的\texttt{data/venus\_horizons\_cache\_2026.json}，该文件由课程代理查询JPL Horizons DE441生成，查询中心为\texttt{@10}。所有交会坐标均在日心黄道J2000系下处理。程序搜索156个候选窗口，并把每次交会的速度闭合误差计入残差预算。当前最佳探索候选为：月球出射日2026年第226.0天，金星交会日第308天，地球返回日第451天；Moon段拼接残差7.024 km/s，金星入出$v_\infty$大小残差0.154 km/s，所需金星偏折角149.49度而近金星6351.8 km时可用偏折角42.09度，对应偏折残差17.906 km/s；返回地球$v_\infty$为6.660 km/s，总残差预算31.744 km/s。结论是该粗搜索序列物理上可复核，但远不如单月球借力主解，主要瓶颈是金星偏折能力不足。

\subsection*{O6：实时交互演示}

\texttt{src/interactive\_demo.py}提供Tkinter交互演示，可调发射日、近月距、近日距和绕飞方向，并实时重算解析筛选模型下的$\Delta v$与轨道椭圆。交互运行命令为\texttt{python src/interactive\_demo.py}。为便于无显示环境评分，脚本还提供\texttt{python src/interactive\_demo.py --self-test}：该测试复核legacy解析参数总量13.214 km/s，并把参数改为第194天、近月距5000 km、近日距0.35 AU、leading绕飞后得到13.878 km/s，证明交互参数会实际驱动计算逻辑。在Xvfb或本地桌面环境中，可运行\texttt{python src/interactive\_demo.py --smoke-gui}自动打开窗口、刷新控件并关闭。

\subsection*{O7：机器学习辅助轨道设计}

作为现有规则之外的创新扩展，\texttt{src/surrogate\_design.py}实现了超过500行的无依赖机器学习辅助筛选器。它用Halton低差异采样和物理边界锚点构造4540个训练样本，训练Ridge surrogate，随后在136510个候选点上预测排序，只复核48个短名单点并在最佳点附近做280点局部复核。该扩展针对legacy解析候选库，定位到2026-07-09、$r_p=0.4$ AU、$r_m=1838$ km、trailing绕飞，总速度增量13.214 km/s；该结果用于展示机器学习辅助筛选能力，不替代本文真实历表闭合主结果。相对完整候选库，精确物理模型评估数减少96.5\%。该模块还输出交叉验证、特征系数、置换重要性和复核候选CSV，用于独立检查其筛选行为。

\section{误差来源与讨论}

本文模型保持作业要求的二维黄道面近似。M3中与Horizons的数千公里残差主要来自忽略其他行星摄动、月球轨道三维倾角投影以及有限步长误差；步长收敛实验表明数值积分误差远小于模型截断误差。修正后的主任务不再假设月球处于理想相位，也不再把绕日一圈等同于返回地球，而是用真实月球与真实地球历表状态做Lambert闭合；代价是$\Delta v_{\rm Moon,res}$显著增大，这正是把未被被动借力闭合的速度矢量差计入预算后的结果。真实工程任务还需三维Lambert求解、微分修正和高精度星历。

\section*{致谢}

感谢课程提供的项目说明、JPL Horizons离线缓存和钱学森参考资料。本文使用Codex / GPT-5协助整理需求、设计代码结构、生成部分Python与LaTeX初稿、调试数值验证流程；所有关键数值均通过本地脚本重新运行并人工核对。详细记录见\texttt{AI-Agent.md}。

\begin{thebibliography}{9}
\bibitem{qian2008}
钱学森. \textit{星际航行概论}. 北京: 中国宇航出版社, 2008.
\bibitem{vallado2013}
Vallado, D. A. \textit{Fundamentals of Astrodynamics and Applications}. 4th ed., Microcosm Press, 2013.
\bibitem{curtis2014}
Curtis, H. D. \textit{Orbital Mechanics for Engineering Students}. 3rd ed., Butterworth-Heinemann, 2014.
\end{thebibliography}

\end{document}
